\section{Compact Appendix A: kernel closure floor}
\label{app:compactA}

This appendix is a compact foundational floor rather than an audit-edition battlefield. Its job is limited and exact: to show why the target manuscript may begin from a kernel of inferential governance rather than from bare derivation-trace classification, and why that kernel already attaches at Level~2 inferential governance rather than only at the stronger Level~3 closure-stable certification layer. The appendix is eliminative. It shows why weaker verdict-regimes fail, why inferential governance is the first regime strong enough for same-domain closure-sensitive reasoning, and why the kernel is forced, mutually closed, and not faithfully generated from below within the same meaningful regime.

\subsection*{A.1. Trace, governance, and certification}

\begin{definition}[Formal calculus and derivation-check]
A formal calculus is a presentation-level object consisting of strings, formation rules, axiom schemata or primitive judgements, and step rules generating finite derivation traces. A derivation-check is the finite verification problem ``does this trace conform to the calculus?'' and yields only presentation-relative rule-conformity.
\end{definition}

\begin{definition}[Displayed package and package-preserving reformulation]
A displayed package is a quintuple $P=(T,B,R,S,I)$ consisting of a displayed target token $T$, displayed support basis $B$, displayed rule/dependency specification $R$, scope tag $S$, and presentation-identity marker $I$. A package-preserving reformulation changes only skin-level presentational features while leaving $P$ fixed. Target drift, support drift, scope drift, identity drift, and verdict drift are any untracked changes in those coordinates or in the licensed / unlicensed verdict attached to the same displayed package.
\end{definition}

\begin{definition}[Inferential governance and closure-stable certification]
A derivation-trace is bare when it is assessed only for rule-conformity. It becomes inferentially governed when, through a claim-bridge, its verdict is used as a licit proposition-bearing basis for assertion, rejection, critique, provisional reuse, or downstream inferential extension. It is closure-stably certified when, in addition, that licensed status is fixed at act-time and remains stable under lawful redescription and retained-lineage continuation.
\end{definition}

\begin{theorem}[Syntax-standing separation]
Rule-conformity by itself does not yet realize inferential governance. A syntactic verdict becomes governance-bearing only if it stably governs the licit downstream uptake of the same displayed package.
\end{theorem}

\begin{proof}
A derivation-check determines only that a trace conforms to some displayed calculus. It does not yet fix whether the checked trace is being used as a licit claim-bearing contribution, whether the same displayed package remains fixed under package-preserving reformulation, or whether verdict differences make any downstream inferential difference. Those additional burdens belong to inferential governance, not to syntax alone.
\end{proof}

\begin{theorem}[Governance realization criterion]
A verdict mechanism realizes inferential governance for a displayed package if and only if all of the following hold: verdict stability under package-preserving reformulation; explicit reset or re-admission under target / support / scope / identity drift; downstream sensitivity of verdict differences; and no same-package repair by role-shift or silent redescription.
\end{theorem}

\begin{theorem}[Explicit non-vacuity families]
Not every syntactic gate qualifies for inferential governance. In particular, formatting-sensitive and support-order-sensitive syntactic gates fail the governance realization criterion.
\end{theorem}

\begin{proof}
A formatting-sensitive gate that flips verdict when the same trace is re-presented with a boldface change fails package-preserving verdict stability. A support-order-sensitive gate that privileges one displayed order of the same support basis likewise fails package-stability or permits untracked support drift. These families are mathematically explicit and show that ``syntactic gate'' is strictly weaker than inferential governance.
\end{proof}

\subsection*{A.2. Weaker uptake exhaustion}

\begin{proposition}[Weaker-uptake decomposition]
Any verdict-regime weaker than inferential governance weakens at least one of exactly four constraints: package-stable verdicts, anti-drift, anti-repair, or downstream sensitivity.
\end{proposition}

\begin{theorem}[Weaker-uptake exhaustion]
The minimal structurally distinct weaker competitors are exactly: trace-only uptake, drift-tolerant uptake, revision-tolerant same-package uptake, and support-flexible uptake. No fifth minimal weaker competitor remains.
\end{theorem}

\begin{proof}
By the governance realization criterion, any weaker regime must fail at least one of the four constitutive constraints. Grouping failures by which constraint is weakened yields the four listed competitor families. Any purported fifth regime either reduces to one of them or simply satisfies the full criterion and is not weaker after all.
\end{proof}

\begin{theorem}[Characterization of inferential governance]
Inferential governance is the first verdict-regime strong enough to support proposition-bearing same-domain reasoning with stable package identity, closure-relevant downstream uptake, and no same-package repair.
\end{theorem}

\subsection*{A.3. Kernel attachment and closure}

\begin{definition}[Non-degenerate reasoning]
A reasoning act is non-degenerate exactly when determinate reference, standing, and irreversibility hold together. Admissibility names the governance boundary under which those conditions are jointly enforceable on the same act.
\end{definition}

\begin{theorem}[Minimal assessability boundary]
For inferentially governed proposition-bearing reasoning acts, reference, standing, and irreversibility are jointly necessary and minimal. Remove any one of them and the target phenomenon collapses.
\end{theorem}

\begin{proof}
Without reference, no determinate proposition remains. Without standing, no act qualifies as a licit reasoning step. Without irreversibility, there is no stable difference between valid continuation and retrospective rescue. Any additional independent governance requirement would over-govern the target beyond what assessability itself requires.
\end{proof}

\begin{theorem}[Kernel attachment at Level~2 inferential governance]
Once a verdict mechanism realizes inferential governance in the present sense, admissibility, standing, reference, and irreversibility are already in force. The kernel therefore first attaches at Level~2 inferential governance rather than only at the stronger Level~3 closure-stable certification layer.
\end{theorem}

\begin{proof}
Inferential governance already requires stable package identity, proposition-bearing use, licensed uptake, and anti-repair. Those are exactly the local work of reference, standing, admissibility, and irreversibility. Closure-stable certification strengthens this by retained-lineage hardening, but it does not relocate the first attachment point of the kernel.
\end{proof}

\begin{theorem}[Kernel non-derivability and no faithful lower generator]
Within the same meaningful regime, no faithful lower generator derives admissibility, standing, reference, or irreversibility from below. Any such attempt either presupposes the kernel, changes the target regime, or collapses into bookkeeping.
\end{theorem}

\begin{theorem}[Mutual closure and minimality of the kernel]
The kernel is mutually closed and minimal for the target phenomenon. No proper subset suffices for inferentially governed proposition-bearing same-domain reasoning, and no additional independent foundational condition survives beside it in the same regime.
\end{theorem}

\begin{theorem}[Level~2 / Level~3 separation]
Inferential governance and closure-stable certification are distinct. Level~2 supplies licensed proposition-bearing use. Level~3 adds closure-stable certification and retained-lineage hardening. The latter is stronger, but it does not define the former.
\end{theorem}

\begin{remark}[Kernel downstream forcing note]
Once the compact Appendix~A floor is fixed, the interface-forcing, quotient, and coherent-reasoning layers are treated as forced consequences relative to that kernel. The coherent-reasoning layer then closes the hardest hidden-path attack: coherence is pointwise standing-preservation along legal same-scope trajectories; no deferred coherence exists; terminal standing certifies coherence; all coherence-relevant distinctions collapse to the standing quotient and tensor invariants; and no creation, amplification, repair, or recovery of standing survives inside the regime. Consequently, any same-domain G\"odel threat that would matter to inferential life must survive into standing classification or coherent trajectory structure. If it survives nowhere there, it is closure-inert inside the regime.
\end{remark}
